\documentclass[11pt]{article}
\usepackage[margin=1in]{geometry}
\usepackage{amsmath,amssymb,amsthm}
\usepackage{xurl}
\usepackage{hyperref}
\sloppy
\let\oldus\_ \renewcommand{\_}{\oldus\allowbreak}
\newtheorem{theorem}{Theorem}
\newtheorem{lemma}[theorem]{Lemma}
\theoremstyle{definition}
\newtheorem{definition}[theorem]{Definition}
\newtheorem{remark}[theorem]{Remark}
\newcommand{\tr}{\operatorname{tr}}
\newcommand{\Fix}{\operatorname{Fix}}
\newcommand{\Z}{\mathbb{Z}}
\newcommand{\R}{\mathbb{R}}
\newcommand{\C}{\mathbb{C}}
\newcommand{\T}{\mathbb{T}}

\title{Disproof of conjecture 00000001485:\\ the cat-map suspension has zeta radius at most $1/2$, not $1$}
\author{Jackmeson1}
\date{2026-10-05}

\begin{document}
\maketitle

\section{The conjecture and the clause refuted}

\textbf{Statement (English, verbatim).} Definition: The zeta function of a suspension of an invertible
map is the generating function $\zeta(z)=\prod(1-z^p)^{-1}$ of its periodic counts. Conjecture: For
every hyperbolic suspension, $\zeta$ has radius of convergence 1, and its singular set on the unit
circle is exactly the set of roots of unity (a complete characterization of periodic structure); no
hyperbolic suspension has radius of convergence $>1$.

The conjecture is a conjunction. \textbf{We refute only its first clause}, ``for every hyperbolic
suspension, $\zeta$ has radius of convergence $1$''. That suffices to refute the conjunction. The
witness is the suspension flow of Arnold's cat map on the $2$-torus. For it we prove that the
radius of convergence of $\zeta$ is at most $1/2$. Its true value is $(3-\sqrt5)/2$, but only the
bound is proved in Lean. The clause about the singular set is not formalized (see
Remark~\ref{rem:scope}). Neither is the clause ``no hyperbolic suspension has radius $>1$'', which
the example satisfies.

\section{Reading and definitions (as formalized)}

All objects below are Lean/Mathlib definitions in \texttt{Conjecture1485/Basic.lean}, namespace
\texttt{C1485}. Here $\mathbb N=\{0,1,2,\dots\}$.

\begin{definition}[Torus and toral maps]
$\T^2=\R^2/\Z^2$ is the quotient of the additive group $\R^2=$ \texttt{Fin 2 -> R} by the lattice
\texttt{Lat}, the image of $\Z^2$ (\texttt{Torus}). An integer matrix $M$ induces the group
endomorphism \texttt{tmap M}, $x\bmod\Z^2\mapsto Mx\bmod\Z^2$.
\end{definition}

\begin{definition}[Hyperbolic toral automorphism]
\texttt{IsHyperbolic M} holds when $\det M$ is a unit of $\Z$ (so $M\in GL(2,\Z)$) and every complex
eigenvalue $\mu$ of $M$ (\texttt{Module.End.HasEigenvalue} of \texttt{Matrix.toLin'}) has
$|\mu|\neq 1$.
\end{definition}

\begin{definition}[Cat map]
$A=\begin{pmatrix}2&1\\1&1\end{pmatrix}$, with inverse $\begin{pmatrix}1&-1\\-1&2\end{pmatrix}$. The
cat map \texttt{cat} is \texttt{tmap A}, packaged as an additive automorphism of $\T^2$ with inverse
\texttt{tmap Ainv}. So it is an invertible map.
\end{definition}

\begin{definition}[Suspension and suspension flow]
For a bijection $e$ of a set $\alpha$, \texttt{Suspension e} is the quotient of $\alpha\times\R$ by
$(x,t)\sim(e^kx,\,t-k)$ for $k\in\Z$. This is the equivalence relation generated by
$(x,t+1)\sim(e\,x,t)$, i.e.\ the mapping torus with unit roof. The flow is
$\varphi_s[x,t]=[x,t+s]$ (\texttt{flow}).
\end{definition}

\begin{definition}[Periodic orbits, periodic counts, $\zeta$, radius]
For $f:\alpha\to\alpha$, a \emph{periodic orbit} (\texttt{PerOrbit f}) is a set
$\{f^kx:k\in\mathbb N\}$ with $x$ a periodic point. Its period $|\tau|$ (\texttt{PerOrbit.len}) is its
number of points, and it equals the least period of $x$ (\texttt{PerOrbit.len\_eq}). The periodic
counts are $N_n=\#\Fix(f^n)$. The zeta function is the infinite product over all periodic orbits,
\[
\zeta_f(z)=\prod_{\tau}(1-z^{|\tau|})^{-1}\in\R[[z]]
\qquad(\texttt{zeta f := tprod (fun tau => fac tau.len)}),
\]
with the Euler factor $\texttt{fac }p=(1-z^p)^{-1}$ (Mathlib's power-series inverse). The product is
taken in Mathlib's coefficientwise topology on $\R[[z]]$ (\texttt{PowerSeries.WithPiTopology}). Each
orbit contributes its own factor, so the periodic counts enter with multiplicity. The radius of
convergence \texttt{zetaRadius f} is Mathlib's \texttt{FormalMultilinearSeries.radius} of
$\sum_n \zeta_n z^n$ (\texttt{ofScalars} over $\C$), i.e.\
$\sup\{r\ge0: (|\zeta_n|r^n)_n \text{ bounded}\}$.
\end{definition}

\begin{definition}[The refuted clause]
\texttt{Clause1}: for every integer matrix $M$ with \texttt{IsHyperbolic M},
$\texttt{zetaRadius (tmap M)}=1$. This is the first clause, restricted to suspensions of hyperbolic
automorphisms of $\T^2$. That is a subclass of all hyperbolic suspensions, so refuting it there
refutes the ``for every'' claim.
\end{definition}

\section{Proof}

\begin{lemma}[\texttt{A\_hyperbolic}]
$A$ is hyperbolic.
\end{lemma}
\begin{proof}
$\det A=1$. An eigenvalue $\mu$ is a root of the characteristic polynomial $z^2-3z+1$
(\texttt{charpoly\_fin\_two}). Write $\mu=a+bi$ with $a^2+b^2=1$. The real and imaginary parts give
$a^2-b^2-3a+1=0$ and $b(2a-3)=0$, hence $a(2a-3)=0$. If $a=0$ then $b=0$, contradicting
$a^2+b^2=1$. If $a=3/2$ then $a^2>1$, again a contradiction.
\end{proof}

\begin{lemma}[\texttt{card\_ker\_tmap}]
If $\det M\neq0$, the kernel of \texttt{tmap M} is finite with exactly $|\det M|$ elements.
\end{lemma}
\begin{proof}
Let $\psi:\Z^2\to\T^2$, $\psi(v)=M^{-1}v\bmod\Z^2$. Then $\ker\psi=M\Z^2$, and the range of $\psi$ is
$\ker(\texttt{tmap }M)$: if $Mx\in\Z^2$ then $x=M^{-1}(Mx)$. So $\ker(\texttt{tmap }M)\cong\Z^2/M\Z^2$.
Mathlib's \texttt{Submodule.natAbs\_det\_equiv}, applied to the injective map $v\mapsto Mv$, gives
$\#(\Z^2/M\Z^2)=|\det M|$.
\end{proof}

\begin{lemma}[\texttt{natAbs\_det\_ge}, \texttt{fix\_cat}]
For $n\ge1$, $\Fix(\mathrm{cat}^n)$ is finite and $N_n=|\det(A^n-I)|\ge 2^n-1$.
\end{lemma}
\begin{proof}
$\mathrm{cat}^n=\texttt{tmap}(A^n)$ (\texttt{cat\_iterate}), so $\Fix(\mathrm{cat}^n)=\ker\texttt{tmap}(A^n-I)$.
For $2\times2$ matrices, $\det(B-I)=\det B-\tr B+1$, and $\det A^n=1$, so $\det(A^n-I)=2-t_n$ with
$t_n=\tr A^n$. From $A^2=3A-I$ we get $t_{n+2}=3t_{n+1}-t_n$ (\texttt{trace\_rec}), with $t_0=2$ and
$t_1=3$. By induction, $2t_{n+1}\le t_{n+2}$ and $t_{n+1}\ge2^{n+1}+1$ (\texttt{trace\_growth}). Hence
$|\det(A^n-I)|=t_n-2\ge2^n-1>0$, and the previous lemma applies.
\end{proof}

\begin{lemma}[\texttt{flow\_periodic\_iff}]
For $s>0$: $\varphi_s[x,t]=[x,t]$ iff $s=n$ for some integer $n>0$ with $e^n x=x$.
\end{lemma}
\begin{proof}
$[x,t+s]=[x,t]$ iff there is $k\in\Z$ with $x=e^kx$ and $t=t+s-k$. That is, $s=k$, and $k>0$
because $s>0$.
\end{proof}
So the closed orbits of the suspension flow are exactly the flow lines through the periodic points
of $e$, and the least period of the flow line through $[x,t]$ is the least period of $x$ under $e$.
The periodic orbits and periods of the suspension are therefore those of the cat map, which is how
$\zeta$ is built.

\begin{lemma}[\texttt{finite\_orbits}, \texttt{hasProd\_zeta}]
If every $\Fix(f^m)$ with $m\ge1$ is finite, then for each $k$ there are finitely many orbits with
$|\tau|\le k$ (call this set $F_k$). The product defining $\zeta_f$ converges (\texttt{HasProd}), and
$\zeta_k$ equals the $z^k$-coefficient of the finite product $\prod_{\tau\in F_k}(1-z^{|\tau|})^{-1}$.
\end{lemma}
\begin{proof}
For $p>k$, $(1-z^p)^{-1}-1=z^p(1-z^p)^{-1}$ is divisible by $z^{k+1}$, and a product of series that are
$\equiv1 \pmod{z^{k+1}}$ is again $\equiv1$. Hence for every finite set $S\supseteq F_k$ of orbits, the
$z^k$-coefficient of $\prod_{\tau\in S}$ equals that of $\prod_{\tau\in F_k}$. So the partial products
converge coefficientwise, and \texttt{tprod} is their limit.
\end{proof}

\begin{lemma}[\texttt{fac\_eq}, \texttt{prod\_ge}, \texttt{coeff\_zeta\_ge}]
For $k\ge1$, $\zeta_k\ge\#\{\tau\in F_k: |\tau|\text{ divides }k\}$.
\end{lemma}
\begin{proof}
$(1-z^p)^{-1}=\sum_{p\mid j}z^j$ for $p\ge1$, so each factor $G_\tau$ has $G_\tau-1$ with nonnegative
coefficients. By induction over a finite set, $\prod G_\tau-1-\sum(G_\tau-1)$ and $\prod G_\tau-1$
have nonnegative coefficients. Taking the $z^k$-coefficient gives
$\zeta_k\ge\sum_{\tau\in F_k}[\,|\tau|\mid k\,]$.
\end{proof}

\begin{lemma}[\texttt{fix\_le\_orbits}]
For $k\ge1$, $\#\Fix(f^k)\le k\cdot\#\{\tau\in F_k:|\tau|\mid k\}$.
\end{lemma}
\begin{proof}
Each $x\in\Fix(f^k)$ is periodic, and its orbit $\tau$ has $|\tau|\mid k$, so $\tau\in F_k$. The orbit
map $x\mapsto\tau$ has fibres contained in $\{f^jx_0:j<k\}$, which have at most $k$ points.
\end{proof}

\begin{lemma}[\texttt{radius\_le\_half}]
If $n c_n\ge2^n-1$ for all $n\ge1$, the series $\sum c_n z^n$ has radius $\le1/2$.
\end{lemma}
\begin{proof}
Suppose the radius exceeds $1/2$. Pick $r$ with $1/2<r<\min(\text{radius},1)$. Then
$|c_n|r^n\le C$ for all $n$ (\texttt{norm\_mul\_pow\_le\_of\_lt\_radius}), so
$(2r)^n=2^nr^n\le(nc_n+1)r^n\le nC+1$. This contradicts $n/(2r)^n\to0$, since $2r>1$.
\end{proof}

\begin{theorem}[\texttt{conjecture\_00000001485\_false}]
$A$ is hyperbolic. For all $n\ge1$, $\#\Fix(\mathrm{cat}^n)=|\det(A^n-I)|$. The closed orbits of the
suspension flow of \texttt{cat} are characterized as in \texttt{flow\_periodic\_iff}. The product
$\prod_\tau(1-z^{|\tau|})^{-1}$ converges to $\zeta_{\mathrm{cat}}$. And
$\texttt{zetaRadius cat}\le1/2$, hence $\neg\,\texttt{Clause1}$.
\end{theorem}
\begin{proof}
Combining the lemmas, for $n\ge1$: $2^n-1\le N_n\le n\cdot\#\{\tau\in F_n:|\tau|\mid n\}\le n\,\zeta_n$.
So \texttt{radius\_le\_half} applies. If \texttt{Clause1} held, then applying it to $M=A$ would give
radius $1\le1/2$, which is false.
\end{proof}

\begin{remark}[Scope and readings]\label{rem:scope}
(i) Only clause 1 is refuted. The singular-set clause and the clause ``no radius $>1$'' are not
addressed in Lean. As a prose remark only, not formalized: clause 2 cannot hold for any power series
of radius $1$. The set of points on the circle of convergence where the series has no analytic
continuation is closed, while the set of roots of unity is not closed.
(ii) ``Hyperbolic suspension'' is instantiated by the subclass of suspensions of hyperbolic
automorphisms of $\T^2$. General Anosov flows are not formalized.
(iii) $\prod(1-z^p)^{-1}$ is read as a product over periodic orbits, so the periodic counts appear
with multiplicity. This is equivalent to $\exp(\sum_n N_n z^n/n)$; that equivalence is not used or
formalized. The alternative reading, a product over the \emph{set} of periods without multiplicity,
is not refuted. It does not encode any periodic counts.
(iv) Conventions: Mathlib's \texttt{Function.periodicPts} requires a positive period, so every
$|\tau|\ge1$ (\texttt{PerOrbit.len\_pos}) and every Euler factor is a genuine inverse.
\end{remark}

\section{Relation to the earlier submission}

The earlier submission is PR \#262 by orionsheep, closed without merging. The reviewer wrote (Unicode
symbols transliterated): ``conjecture\_refuted is a conjunction of decidable trivia (four periodic
counts, Z[$\sqrt5$] identities, sign flags) that never mentions $\zeta$, the suspension, or a radius
of convergence. The cat-map counterexample --- valid in prose --- is entirely unformalized, and only
finitely many periodic counts are certified where the radius of convergence needs all of them.''

This submission uses the same counterexample, and formalizes everything that the earlier Lean lacked:
\begin{itemize}
\item the torus $\R^2/\Z^2$, with the cat map as an actual automorphism of it, and a proof of
hyperbolicity;
\item the periodic counts $\#\Fix(\mathrm{cat}^n)=|\det(A^n-I)|$ for \emph{every} $n\ge1$;
\item the suspension flow and its closed orbits;
\item $\zeta$ as the convergent infinite product over all periodic orbits;
\item Mathlib's radius of convergence of $\zeta$, bounded using all coefficients.
\end{itemize}

\section{Lean formalization and axiom audit}

Package: \texttt{lean/Conjecture1485/Basic.lean} (445 lines, \texttt{import Mathlib} only; Lean
v4.33.1, Mathlib v4.33.1), identical to the file built with \texttt{lake build}. The main theorem is
\texttt{C1485.conjecture\_00000001485\_false}. The lemmas are named in the sections above.
\texttt{\#print axioms C1485.conjecture\_00000001485\_false} reports
\texttt{[propext, Classical.choice, Quot.sound]}. There is no \texttt{sorry} and no
\texttt{native\_decide}. \texttt{decide} is used only for $2\times2$ integer matrix identities
($AA^{-1}=I$, $A^2=3A-I$, $\det A=1$, $\tr A^0=2$, $\tr A=3$).

\section*{Sources}
\begin{itemize}
\item Earlier submission and review: \url{https://github.com/The-Last-Math-Competition/The-Last-Math-Competition/pull/262}.
\item Mathlib v4.33.1 (definitions and the theorems named above, e.g.\
\texttt{Mathlib/LinearAlgebra/FreeModule/Finite/CardQuotient.lean}).
\end{itemize}

\end{document}
